\scanpage{79}
% archive sheet number: 78
% L1-2, bibliographic notes crossed together by a wavy stroke
\deleted{F. \unsure{Peterson|Petterson} Functional Cohomology operations. Some non-embedding Problems.}

\deleted{Moore The double suspension and $p$-primary components of homotopy gps of spheres.}
% L3-5; cancellation strokes through paired factors
\[
u^{-1}\alpha\,\deleted{Tu\,Tu^{-1}}\,T\alpha\,\deleted{T^2u\,T^2u^{-1}}\,T^2\alpha\,\deleted{T^3u}\cdots\deleted{T^{n-1}u^{-1}}\,T^{n-1}\alpha\,T^nu=1
\]
i.e.
\[
u^{-1}(\alpha\,T\alpha\cdots T^{n-1}\alpha)T^nu=1
\]
d'où, comme $(\quad)=1$, $u^{-1}T^nu=1$ soit $T^nu=u$.
% L6-9
\uline{Lemme 2} Si un élément de $H^1(\mathbb Z_n,G_n)$ est représenté par un $\alpha\in G_n$ tel que $\alpha\,T\alpha\cdots T^{n-1}\alpha=1$, alors son image dans $H^1(\mathbb Z_m,G_m)$ (si $n\mid m$) est de $=$ représentée par $\alpha\in G_n\subset G_m$.

\unsure{Trivial|Triviale}
% L10
Posons
\[
{\unsure{\mathscr H|\mathcal H|\mathbb H}}^{\unsure{1}}(T,G)=\underset{\substack{\unsure{n|\longrightarrow}\\\text{suivant l'ordre}\\\text{filtrant arithmétique}}}{\lim}H^1(\mathbb Z_n,G_n).
\]
% L11-14
On voit donc que $H^1(T,G)$ est l'ens des classes d'objets $\alpha\in G$ tels qu'il existe un $n$ avec $\alpha\,T\alpha\cdots T^{n-1}\alpha=1$, \deleted{\unsure{et}} en écrivant $\beta\sim\alpha$ s'il existe $u\in G$ tel que $\beta=u^{-1}\alpha\,Tu$.
% L15-17
\uline{Lemme 3} Si \unsure{Pour} \deleted{\unsure{tout}\ \unsure{avec|groupe}\ $G$} l'application $x\unsure{\longrightarrow|\text{---}}x^{-1}Tx$ de $G$ dans $G$ est surjective, on a $H^1(T,G)=\{e\}$.

Facile sur lemmes 1 et 2
% L18-23
\uline{Lemme 4} Soit
\diagram{\begin{tikzcd}[ampersand replacement=\&]
e\arrow[r,dash]\&G'\arrow[r,dash]\&G\arrow[r,dash]\&G''\arrow[r,dash]\&e
\end{tikzcd}}
une suite exacte de $\mathbb Z$-groupes ($T$ opère dans $G$ et conserve $G'$).
Supposons que pour $\forall$ automorphisme \inserted{intérieur $\sigma$} de $G$, \deleted{\unsure{on a}} \deleted{\unsure{pour chaque hom intérieur de $G$ l'application|pour chaque automorphisme intérieur de $G$ l'application}} \inserted{l'image de $G$ par l'application} $x\unsure{\longrightarrow|\text{---}}x^{-1}\sigma Tx$ \deleted{de $G'$ dans $G'$ soit surjective} \inserted{contienne $G'$}. Alors
\[
H^1(T,G)\unsure{\longrightarrow|\text{---}}H^1(T,G'')\quad\text{est \uline{injectif}}
\]
% L24-25
Soient en effet $\alpha,\beta\in G$, tels que $\alpha\,T\alpha\cdots T^{n-1}\alpha=1$ et de $=$ $\beta\,T\beta\cdots T^{n-1}\beta=1$, et
